% ═════════════════════════════════════════════════════════════════════════════
\section{Propositional Logic}
% ═════════════════════════════════════════════════════════════════════════════

These notes develop the core tools of propositional logic: truth tables,
algebraic proof of logical equivalence, validity of arguments, and the
negation of quantified statements.

% ═════════════════════════════════════════════════════════════════════════════
\section{Truth Tables and Propositional Form}
% ═════════════════════════════════════════════════════════════════════════════

A \defn{truth table} lists the truth value of a compound proposition for every
possible assignment of truth values to its constituent variables. With $n$
propositional variables, there are $2^n$ rows.

\subsection{A worked example}

Consider the proposition
\[
    X = (p \to q) \lor (p \land r).
\]
Since $X$ depends on three variables, there are $2^3 = 8$ rows. We use the
equivalence $p \to q \equiv {\sim}p \lor q$ to expand the implication before
constructing the table.

\[
\begin{array}{|c|c|c|c|c|c|c|c|c|}
    \hline
    p & q & r & {\sim}p & p \land r & {\sim}p \lor q & p \to q
      & (p \to q) \lor (p \land r) & X \\
    \hline
    0 & 0 & 0 & 1 & 0 & 1 & 1 & 1 & 1 \\
    0 & 0 & 1 & 1 & 0 & 1 & 1 & 1 & 1 \\
    0 & 1 & 0 & 1 & 0 & 1 & 1 & 1 & 1 \\
    0 & 1 & 1 & 1 & 0 & 1 & 1 & 1 & 1 \\
    1 & 0 & 0 & 0 & 0 & 0 & 0 & 0 & 0 \\
    1 & 0 & 1 & 0 & 1 & 0 & 0 & 1 & 1 \\
    1 & 1 & 0 & 0 & 0 & 1 & 1 & 1 & 1 \\
    1 & 1 & 1 & 0 & 1 & 1 & 1 & 1 & 1 \\
    \hline
\end{array}
\]

\subsection{Consequences of the truth table}

\begin{proposition}
$X$ is not a tautology.
\end{proposition}
\begin{proof}
Row 5 of the table (where $p = 1$, $q = 0$, $r = 0$) gives $X = 0$.
Since $X$ is false for at least one assignment, it is not a tautology.
\end{proof}

\begin{proposition}
If $p$ is true, $X$ can be false.
\end{proposition}
\begin{proof}
When $p = 1$, $q = 0$, and $r = 0$, the table gives $X = 0$.
\end{proof}

\begin{proposition}
If $q$ is true, $X$ cannot be false.
\end{proposition}
\begin{proof}
The unique row where $X = 0$ is row 5, in which $q = 0$. Hence no assignment
with $q = 1$ makes $X$ false.
\end{proof}

% ═════════════════════════════════════════════════════════════════════════════
\section{Proving Logical Equivalence}
% ═════════════════════════════════════════════════════════════════════════════

Two propositions are \defn{logically equivalent} if they share the same truth
value under every assignment of truth values to their variables. Rather than
constructing a joint truth table, equivalence can often be established by
applying the standard algebraic laws of propositional logic (Implication,
De Morgan, Commutativity, Distribution, etc.) in sequence.

\begin{proposition}
The propositions $p \to (q \lor {\sim}r)$ and ${\sim}(p \land r \land {\sim}q)$
are logically equivalent.
\end{proposition}

\begin{proof}
\begin{align*}
    p \to (q \lor {\sim}r)
    &\equiv {\sim}p \lor (q \lor {\sim}r)
    \tag{Implication: $\phi\to\psi\equiv{\sim}\phi\lor\psi$}\\
    &\equiv {\sim}p \lor {\sim}r \lor q
    \tag{Commutativity and Associativity}\\
    &\equiv {\sim}(p \land r \land {\sim}q).
    \tag{De Morgan: ${\sim}(A\land B\land C)\equiv{\sim}A\lor{\sim}B\lor{\sim}C$}
\end{align*}
Hence $p \to (q \lor {\sim}r) \equiv {\sim}(p \land r \land {\sim}q)$.
\end{proof}

% ═════════════════════════════════════════════════════════════════════════════
\section{Argument Validity via Proof by Contradiction}
% ═════════════════════════════════════════════════════════════════════════════

An argument with premises $P_1, \ldots, P_n$ and conclusion $C$ is
\defn{valid} if and only if it is impossible for every premise to be true
while the conclusion is false. Equivalently: assume all premises true and
$C$ false, then derive a contradiction.

\begin{example}
Determine whether the following argument is valid or invalid.
\[
\begin{array}{l}
    p \lor {\sim}q \\
    {\sim}(p \land r) \\
    r \lor s \\
    \hline
    q \to s
\end{array}
\]
\end{example}

\begin{proof}[Solution]
Assume all premises are true and the conclusion is false:
\begin{alignat}{2}
    &p \lor {\sim}q \equiv 1, \tag{P1}\\
    &{\sim}(p \land r) \equiv 1, \tag{P2}\\
    &r \lor s \equiv 1, \tag{P3}\\
    &q \to s \equiv 0. \tag{C}
\end{alignat}

\noindent\textbf{From (C).}
Using $q \to s \equiv {\sim}q \lor s$, a disjunction is false only when both
disjuncts are false. Therefore $q = 1$ and $s = 0$.

\noindent\textbf{From (P1).}
Since $q = 1$, we have ${\sim}q = 0$, so $p$ must be true to satisfy
$p \lor {\sim}q$. Hence $p = 1$.

\noindent\textbf{From (P2).}
Expanding: ${\sim}(p \land r) \equiv {\sim}p \lor {\sim}r$. Since $p = 1$,
we have ${\sim}p = 0$, so ${\sim}r = 1$, giving $r = 0$.

\noindent\textbf{From (P3).}
With $r = 0$ and $s = 0$, we get $r \lor s = 0 \lor 0 = 0$, contradicting (P3).

The premises cannot all be simultaneously true when the conclusion is false.
Therefore the argument is \textbf{valid}.
\end{proof}

% ═════════════════════════════════════════════════════════════════════════════
\section{Quantifiers and Negation}
% ═════════════════════════════════════════════════════════════════════════════

To negate a quantified statement, each quantifier is switched and the
predicate is negated:
\[
    \lnot\,\forall x\; P(x) \;\equiv\; \exists x\; \lnot P(x),
    \qquad
    \lnot\,\exists x\; P(x) \;\equiv\; \forall x\; \lnot P(x).
\]
When multiple quantifiers appear, they are switched from left to right.

\begin{example}\label{ex:negations}
State the negation of each proposition.
\begin{enumerate}[label=(\alph*)]
    \item $\forall$ $2 \times 2$ matrices $A, B$: $\det(AB) = \det(BA)$.
    \item $\forall$ integers $n \geq 5$, $\exists$ integers $a, b \geq 0$
          such that $n = 3a + 5b$.
    \item $\forall x \geq 4,\; \forall y \geq x,\; \exists z \geq 7$ such
          that $x + y = z$.
\end{enumerate}
\end{example}

\begin{proof}[Negations]
\begin{enumerate}[label=(\alph*)]
    \item $\exists$ $2 \times 2$ matrices $A, B$ such that
          $\det(AB) \neq \det(BA)$.

    \item $\exists$ integer $n \geq 5$ such that $\forall$ integers
          $a, b \geq 0$, $n \neq 3a + 5b$.

    \item $\exists x \geq 4,\; \exists y \geq x$ such that
          $\forall z \geq 7$, $x + y \neq z$.
\end{enumerate}
\end{proof}

\begin{proposition}
Proposition \textup{(c)} in Example~\ref{ex:negations} is true.
\end{proposition}

% CORRECTION: The original argument cited two specific numerical examples
% (x=4,y=4 and x=100,y=101) as justification, but examples do not constitute
% a proof. The correct argument makes the inequality chain explicit.

\begin{proof}
Let $x \geq 4$ and $y \geq x$ be arbitrary real numbers. Since $y \geq x$,
\[
    x + y \geq x + x = 2x \geq 2 \cdot 4 = 8 > 7.
\]
Set $z = x + y$. Then $z \geq 8 \geq 7$, so the condition $z \geq 7$ is
satisfied, and trivially $x + y = z$. Since $x$ and $y$ were arbitrary, the
proposition holds for all valid choices.
\end{proof}
